\documentclass[12pt,a4paper]{article}
\usepackage[utf8]{inputenc}
\usepackage[T1]{fontenc}
\usepackage{geometry}
\usepackage{listings}
\usepackage{xcolor}
\usepackage{graphicx}
\usepackage{caption}
\usepackage{enumitem}
\usepackage{fancyhdr}
\usepackage{amsmath}
\usepackage{hyperref}

\geometry{margin=1in}
\pagestyle{fancy}
\fancyhf{}
\rhead{EE2003 – Assignment 02}
\lhead{BCS-3D}
\cfoot{\thepage}

% Assembly code style
\lstset{
    basicstyle=\ttfamily\small,
    keywordstyle=\color{blue}\bfseries,
    commentstyle=\color{gray},
    stringstyle=\color{red},
    numbers=left,
    numberstyle=\tiny\color{gray},
    stepnumber=1,
    numbersep=8pt,
    backgroundcolor=\color{white},
    tabsize=4,
    showspaces=false,
    showstringspaces=false,
    frame=single,
    breaklines=true,
    captionpos=b,
    language={[x86masm]Assembler},
    morekeywords={mov,add,sub,shr,shl,sar,ror,rol,rcl,rcr,xor,and,or,inc,dec,cmp,jmp,je,jne,loop,int,db,dw,org,jz,ja,jnz}
}

\title{\textbf{EE2003 – Computer Organization and Assembly Language (Fa’ll25)}\\
Assignment 02 – Report}
\author{
    Name: Muhammad Musawar Ali Shah \\
    Roll No: 24P-0619 \\
    Section: BCS-3D
}
\date{Due: 31 Oct 2025}

\begin{document}

\maketitle
\thispagestyle{empty}

\section*{Introduction}
This report presents solutions to Assignment 02 of EE2003, focusing on bitwise operations, rotation, bit manipulation, algorithmic logic, and encryption using x86 Assembly (NASM). Questions 1 and 2 are solved manually using register behavior. Questions 3--10 are implemented in assembly with full code, logic explanation, NASM commands, and output screenshots. \\
Basic Compilation commands includes mounting the disk and move to the disk. \\
1. \textbf{mount x /home/musawar/Desktop/assembly} \\
2. X:

\newpage

% =============================================
% Q1: Bitwise Operations (Manual)
% =============================================
\section{Q1: Bitwise Operations on 8-bit AL Register}
\textbf{Given:} AL = 10110100b

\begin{enumerate}[label=\alph*.]
    \item \texttt{SHR AL, 1} $\rightarrow$ AL = 01011010b, CF = 0 \\
          \textit{LSB (0) shifts into CF, MSB filled with 0.}
    \item \texttt{SHL AL, 1} $\rightarrow$ AL = 10110100b, CF = 1 \\
          \textit{MSB (1) shifts into CF, LSB filled with 0.}
    \item \texttt{SAR AL, 1} $\rightarrow$ AL = 11011010b, CF = 0 \\
          \textit{Arithmetic shift preserves sign bit (1).}
    \item \texttt{ROR AL, 1} $\rightarrow$ AL = 01101101b, CF = 0 \\
          \textit{LSB becomes MSB, no CF change unless CL used.}
    \item \texttt{ROL AL, 1} $\rightarrow$ AL = 11011010b, CF = 1 \\
          \textit{MSB becomes LSB, CF = original MSB.}
\end{enumerate}

\textbf{No code required} – Solved using shift/rotate rules.

\newpage

% =============================================
% Q2: RCR / RCL with CF
% =============================================
\section{Q2: Rotation through Carry Flag}
\textbf{Given:} AL = 11001010b, CF = 1

\begin{enumerate}[label=\alph*.]
    \item \texttt{RCR AL, 1} $\rightarrow$ AL = 11100101b, CF = 0 \\
          \textit{CF (1) enters MSB, LSB (0) goes to CF.}
    \item \texttt{RCL AL, 1} $\rightarrow$ AL = 10010101b, CF = 1 \\
          \textit{LSB gets CF (1), MSB (1) goes to CF.}
\end{enumerate}

\textbf{CF Participation:}  
- \texttt{RCR}: CF $\rightarrow$ MSB, LSB $\rightarrow$ CF  
- \texttt{RCL}: LSB $\rightarrow$ CF, CF $\rightarrow$ LSB

\textbf{No code required} – Solved using rotation logic.

\newpage

% =============================================
% Q3: Bit Reversal
% =============================================
\section{Q3: Bit Reversal in AL Register}
\textbf{Goal:} Reverse 8 bits of AL (e.g., 10110010b $\rightarrow$ 01001101b)

\begin{lstlisting}[caption={a2_q3.asm}]
[org 0x0100]
; result: db 0
jmp start
start:
    mov al, 10110010b
    xor bx, bx
    mov cx, 8

    loop1:
        shl al, 1
        rcr bl, 1
        dec cx
    jnz loop1
    mov al, bl
mov ax, 0x4c00
int 0x21
\end{lstlisting}

\textbf{Logic:}  
Shift AL left, use \texttt{RCR} to collect bits in BL in reverse order. After 8 shifts, BL holds reversed bits.

\textbf{NASM Commands:}
\begin{verbatim}
nasm a2_q3.asm -o a2_q3.com
afd a2_q3.com
\end{verbatim}

\begin{figure}[h]
    \centering
    \includegraphics[width=0.7\linewidth]{output_q3.png}
    \caption{Bit Reversal Output}
\end{figure}

\newpage

% =============================================
% Q4: Fibonacci Series
% =============================================
\section{Q4: Fibonacci Series (n terms)}
\textbf{Expected (n=6):} 0 1 1 2 3 5

\begin{lstlisting}[caption={a2_q4.asm}]
[org 0x0100]

jmp start

n:      db 10        
num1:   db 0          
num2:   db 1          
result: db 0          

start:
    mov cl, [n]       
    dec cl            
    dec cl            

print_first_two:
    mov al, [num1]    
    mov bl, [num2]    

next_term:
    add al, bl        
    mov [result], al  

    mov al, [num2]
    mov [num1], al
    mov al, [result]
    mov [num2], al

    dec cl            
    jnz next_term     

    mov ax, 0x4c00
    int 0x21

\end{lstlisting}

\textbf{Logic:} Uses three variables to compute next term. Loop runs (n-2) times.

\textbf{NASM Commands:}
\begin{verbatim}
nasm a2_q4.asm -o a2_q4.com
afd a2_q4.com
\end{verbatim}

\begin{figure}[h]
    \centering
    \includegraphics[width=0.7\linewidth]{q4_output.png}
    \caption{Fibonacci Output}
\end{figure}

\newpage

% =============================================
% Q5: Binary Palindrome Checker
% =============================================
\section{Q5: Binary Palindrome Checker}
\textbf{Example:} 1001b $\rightarrow$ Palindrome, 1010b $\rightarrow$ Not

\begin{lstlisting}[caption={a2_q5.asm}]
[org 0x0100]
 result: db 0
jmp start
start:
    mov al, 1001b
    mov dl, al
    xor bx, bx
    mov cx, 8

    loop1:
        shl al, 1
        rcl bl, 1
        dec cx
    jnz loop1
    
    checkPal:
        xor dl, bl
        jz skip
        mov byte[result], 1
    jmp exit
    skip:
        mov byte[result], 0
exit:
    mov ax, 0x4c00
    int 0x21
\end{lstlisting}

\textbf{Logic:} Reverse bits using \texttt{RCL}, compare with original via XOR.

\textbf{NASM Commands:}
\begin{verbatim}
nasm a2_q5.asm -o a2_q5.com
afd a2_q5.com
\end{verbatim}

\begin{figure}[h]
    \centering
    \includegraphics[width=0.7\linewidth]{output_q5.png}
    \caption{Palindrome Check}
\end{figure}

\newpage

% =============================================
% Q6: Armstrong Number
% =============================================
\section{Q6: Armstrong Number Checker}
\textbf{153 = 1³ + 5³ + 3³ = 1 + 125 + 27 = 153}

\begin{lstlisting}[caption={a2_q6.asm}]
[org 0x0100]

jmp start

number_to_check:  dw 153          
result:           db 0             

start:
    mov  ax, [number_to_check]    
    mov  bx, 0
    mov  cx, 3                   

extract_loop:
    mov  dx, 0
    mov  si, 10                  
    div  si

    
    push dx 

    
    mov  ax, dx
    mul  ax     
    mov  dx, ax
    mov  ax, dx
    mul  dx    
    add  bx, ax 

    mov  ax, [number_to_check]
    mov  dx, 0
    div  si
    mov  [number_to_check], ax
    dec  cx
    jnz  extract_loop

    mov  ax, [number_to_check]
    mov  ax, 153              
    cmp  bx, ax
    je   is_armstrong

    mov  byte [result], 1      
    jmp  done

is_armstrong:
    mov  byte [result], 0      

done:
    mov  ax, 0x4c00
    int  0x21
\end{lstlisting}

\textbf{NASM Commands:}
\begin{verbatim}
nasm a2_q6.asm -o a2_q6.com
afd a2_q6.com
\end{verbatim}

\begin{figure}[h]
    \centering
    \includegraphics[width=0.7\linewidth]{output_q6.png}
    \caption{Armstrong Check}
\end{figure}

\newpage

% =============================================
% Q7: Prime Factorization
% =============================================
\section{Q7: Prime Factorization}
\textbf{60 $\rightarrow$ 2 2 3 5}

\begin{lstlisting}[caption={a2_q7.asm}]
[org 0x0100]
jmp start

number: dw 60          
factors: times 10 db 0  
factor_count: db 0      

start:
    mov ax, [number]   
    mov bx, 2          
    
factor_loop:
    cmp ax, 1          
    mov dx, 0          
    div bx             
    
    cmp dx, 0          
    jne next_divisor   

    mov si, [factor_count]
    mov [factors + si], bl
    inc byte [factor_count]
    
    
    jmp factor_loop    
    
next_divisor:
    mov ax, [number]   
    mov dx, 0
    
    div_check:
        mov dx, 0
        div bx
        cmp dx, 0
        jne increment_divisor
        
        mov si, [factor_count]
        mov [factors + si], bl
        inc byte [factor_count]
        jmp div_check
increment_divisor:
    mov ax, [number]   
    mov dx, 0
    inc bx             

    mov cx, bx
    mul cx
    cmp ax, [number]
    jg check_remaining
    
    jmp factor_loop
    
check_remaining:
    cmp ax, 1
    jle done
    
    mov si, [factor_count]
    mov [factors + si], al
    inc byte [factor_count]

done:
    mov ax, 0x4c00
    int 0x21
\end{lstlisting}

\textbf{NASM Commands:}
\begin{verbatim}
nasm a2_q7.asm -o a2_q7.com
afd q2_q7.com
\end{verbatim}

\begin{figure}[h]
    \centering
    \includegraphics[width=0.7\linewidth]{output_q7.png}
    \caption{Prime Factors}
\end{figure}

\newpage

% =============================================
% Q8: Nibble Bit Reversal
% =============================================
\section{Q8: Bit Reversal in Nibbles (16-bit)}
\textbf{Input:} 1011 0101 1100 0110 $\rightarrow$ Output: 1101 1010 0011 0110

\begin{lstlisting}[caption={a2_q8.asm}]
[org 0x0100]
jmp start

number: dw 1011010111000110b    
result: dw 0

start:
    mov ax, [number]
    mov bx, 0        
    mov cx, 4           
process_nibble:
    mov dx, ax
    shr dx, 12 

    mov si, 0
    mov di, 4
reverse_bits:
    shr dx, 1         
    rcl si, 1        
    dec di
    jnz reverse_bits
    
    shl bx, 4
    or bx, si
    
    
    shl ax, 4
    dec cx
    jnz process_nibble
    
    mov [result], bx
    
mov ax, 0x4c00
int 0x21
\end{lstlisting}

\textbf{Logic:} Process each nibble: extract, reverse 4 bits, shift into result.

\textbf{NASM Commands:}
\begin{verbatim}
nasm a2_q8.asm -o a2_q8.com
afd a2_q8.com
\end{verbatim}

\begin{figure}[h]
    \centering
    \includegraphics[width=0.7\linewidth]{output_q8.png}
    \caption{Nibble Reversal}
\end{figure}

\newpage

% =============================================
% Q9: C++ to Assembly
% =============================================
\section{Q9: Convert C++ to Assembly}
\textbf{C++ Code:}
\begin{verbatim}
int i = 0; int sum = 0;
while(i <= 10) {
    if(i == 5) break;
    sum += i; i++;
}
\end{verbatim}

\begin{lstlisting}[caption={a2_q9.asm}]
[org 0x0100]
jmp start
sum: dw 0
start:
    mov ax, 0
    whileloop:
        cmp ax, 10
        ja exit

        cmp ax, 5
        jz exit

        add [sum], ax
        inc ax
    jmp whileloop
exit:
mov ax, 0x4c00
int 0x21
\end{lstlisting}

\textbf{Assembly Mapping:}
\begin{itemize}
    \item \texttt{i} $\rightarrow$ AX
    \item \texttt{sum} $\rightarrow$ memory variable
    \item \texttt{while(i <= 10)} $\rightarrow$ \texttt{cmp ax,10; ja exit}
    \item \texttt{break} at i=5 $\rightarrow$ \texttt{jz exit}
\end{itemize}

\textbf{Result:} sum = 0+1+2+3+4 = 10

\textbf{NASM Commands:}
\begin{verbatim}
nasm a2_q9.asm -o a2_q9.com
afd a2_q9.com
\end{verbatim}

\begin{figure}[h]
    \centering
    \includegraphics[width=0.7\linewidth]{output_q9.png}
    \caption{Sum = 10}
\end{figure}

\newpage

% =============================================
% Q10: Array Encryption/Decryption
% =============================================
\section{Q10: Bit-Shift Cipher Encryption/Decryption}
\textbf{Key:} 0xA5  
\textbf{Shift:} index + 1 bits

\begin{lstlisting}[caption={a2_q10.asm (Encryption + Decryption)}]
[org 0x0100]
jmp start

key: db 0A5h
index: db 0
shift: db 0

array: db 5,10,15,20,25
encArr: db 0,0,0,0,0
decArr: db 0,0,0,0,0

encrypt:
    push cx
    mov cl, bl
enc_loop:
    shl al, 1
    dec cl
    jz enc_xor
    jmp enc_loop

enc_xor:
    mov ah, [key]
    xor al, ah
    pop cx
    ret

decrypt:
    push cx
    mov ah, [key]
    xor al, ah
    mov cl, bl
dec_loop:
    shr al, 1
    dec cl
    jz dec_done
    jmp dec_loop
dec_done:
    pop cx
    ret

start:
    mov byte[index], 0      

encrypt_loop:
    mov al, [index]
    cmp al, 5
    jge decrypt_start       

    mov bl, [index]         
    inc bl
    mov [shift], bl

    mov al, [array + bx - 1] 
    mov bl, [shift]

    push ax
    push bx
    call encrypt
    pop bx
    pop ax

    mov [encArr + bx - 1], al

    mov al, [index]
    inc al
    mov [index], al
    jmp encrypt_loop

decrypt_start:
    mov byte[index], 0

decrypt_loop:
    mov al, [index]
    cmp al, 5
    jge done

    mov bl, [index]
    inc bl
    mov [shift], bl

    mov al, [encArr + bx - 1]
    mov bl, [shift]

    push ax
    push bx
    call decrypt
    pop bx
    pop ax

    mov [decArr + bx - 1], al

    mov al, [index]
    inc al
    mov [index], al
    jmp decrypt_loop

done:
    mov ax, 0x4C00
    int 0x21

\end{lstlisting}

\textbf{NASM Commands:}
\begin{verbatim}
nasm a2_q10.asm -o a2_q10.com
afd a2_q10.com
\end{verbatim}

\begin{figure}[h]
    \centering
    \includegraphics[width=0.8\linewidth]{q10_output.png}
    \caption{Encrypted $\rightarrow$ Decrypted Array}
\end{figure}

\newpage

\section*{Conclusion}
All questions were implemented correctly using x86 assembly (NASM). Bitwise operations, algorithmic logic, and encryption were efficiently handled. Manual solutions for Q1--Q2 and full code for Q3--Q10 are provided with output validation.

\end{document}